%--------------------------------------------------------------------%
%
% Berkas utama templat LaTeX.
%
% author Petra Barus, Peb Ruswono Aryan, Faris Rizki Ekananda
%
%--------------------------------------------------------------------%
%
% Berkas ini berisi struktur utama dokumen LaTeX yang akan dibuat.
%
%--------------------------------------------------------------------%

\documentclass[bahasa, 12pt, a4paper, onecolumn, oneside, final]{report}

%-------------------------------------------------------------------%
%
% Konfigurasi dokumen LaTeX untuk laporan tesis IF ITB
%
% @author Petra Novandi
%
%-------------------------------------------------------------------%
%
% Berkas asli berasal dari Steven Lolong
%
%-------------------------------------------------------------------%

% Ukuran kertas
\special{papersize=210mm,297mm}

% Setting margin
\usepackage[top=2cm,bottom=2cm,left=4cm,right=3cm]{geometry}

% Math package
\usepackage{mathptmx}

% 4th Sectioning
\usepackage{titlesec}
\newcommand{\subsubsubsection}[1]{\paragraph{#1}\mbox{}\\}

\titleformat{\subsubsubsection}
{\normalfont\normalsize\bfseries}{\theparagraph}{1em}{}
\titlespacing*{\subsubsubsection}
{0pt}{3.25ex plus 1ex minus .2ex}{1.5ex plus .2ex}

% Judul bahasa Indonesia
\usepackage[bahasa]{babel}

% Format citation
\usepackage[utf8]{inputenc}
\usepackage[style=apa,backend=biber]{biblatex}
\usepackage{longtable}
\usepackage{graphicx}
\usepackage{subfig}
\usepackage{titling}
\usepackage{booktabs}
\usepackage{tabularx}
\usepackage{blindtext}
\usepackage{sectsty}
\usepackage{chngcntr}
\usepackage{etoolbox}
\usepackage{array}
\usepackage[hidelinks]{hyperref}       % Package untuk link di daftar isi. Ubah jadi \usepackage[hidelinks]{hyperref} apabila ingin menghilangkan kotak merah disekitar link
\usepackage{titlesec}       % Package Format judul
\usepackage{titletoc}       % Package Format judul di toc
\usepackage{tocbibind}      % Package untuk masukkan toc, lot, lof ke Daftar Isi
\usepackage{scrwfile}       % Package untuk membuat Daftar Lampiran dari toc
\usepackage{parskip}
\usepackage{afterpage}
\usepackage{relsize}
\usepackage{xcolor, colortbl}
\usepackage{setspace}
\usepackage{listings}

\graphicspath{{resources/}}   % letak direktori penyimpanan gambar

% Setting daftar lampiran
\newcommand*{\lopname}{DAFTAR LAMPIRAN}
\TOCclone[\lopname]{toc}{atoc}
\addtocontents{atoc}{\protect\value{tocdepth}=-1}
\newcommand\listofappendices{
	\cleardoublepage
	\phantomsection
	\listofatoc
	\addcontentsline{toc}{chapter}{\lopname}
}

\newcommand*\savedtocdepth{}
\AtBeginDocument{%
	\edef\savedtocdepth{\the\value{tocdepth}}%
}

\let\originalappendix\appendix
\renewcommand\appendix{%
	\originalappendix
	\cleardoublepage
	\addtocontents{toc}{\protect\value{tocdepth}=-1}%
	\addtocontents{atoc}{\protect\value{tocdepth}=\savedtocdepth}%

	\titlecontents{chapter}
	[0pt]
	{\bfseries}
	{Lampiran \thecontentslabel.\quad}
	{}
	{\hfill\contentspage}

	\titleformat{\chapter}[block]
	{\bfseries}
	{\chaptertitlename\ \thechapter.\quad}{0pt}
	{\bfseries}
}

% Hilangkan titik pada toc
\makeatletter
\renewcommand{\@dotsep}{1}
\makeatother

% Setel title pada chapter-chapter di toc, lof, lot
\titlecontents{chapter}
[0pt]
{\bfseries}
{\MakeUppercase{Bab} \thecontentslabel\quad\uppercase}
{}
{\mdseries\titlerule*[0.35em]{.}\bfseries\contentspage}
\titlecontents{figure}
[0pt]
{}
{Gambar \thecontentslabel.\quad}
{}
{\mdseries\titlerule*[0.35em]{.}\bfseries\contentspage}
\titlecontents{table}
[0pt]
{}
{Tabel \thecontentslabel.\quad}
{}
{\mdseries\titlerule*[0.35em]{.}\bfseries\contentspage}

% Masukin Daftar Pustaka ke toc
\let\originalprintbibliography\printbibliography
\renewcommand\printbibliography{%
	\phantomsection
	\cleardoublepage
	\originalprintbibliography
	\addcontentsline{toc}{chapter}{\bibname}
}

% Line satu setengah spasi
\renewcommand{\baselinestretch}{1.5}

% Setting judul
\chapterfont{\centering \large}
\titleformat{\chapter}[display]
{\Large\centering\bfseries}
{\chaptertitlename\ \thechapter}{0pt}
{\Large\bfseries\uppercase}

% Setting nomor pada subbsubsubbab
\setcounter{secnumdepth}{4}

\makeatletter

\makeatother

% Counter untuk figure dan table.
\counterwithin{figure}{chapter}
\counterwithin{table}{chapter}

% Define blank page
\newcommand*{\blankpage}{\afterpage{\null\newpage}}

% Translate autoref into Indonesian
\renewcommand*{\equationautorefname}{Persamaan}%
\renewcommand*{\footnoteautorefname}{catatan kaki}%
\renewcommand*{\itemautorefname}{item}%
\renewcommand*{\figureautorefname}{Gambar}%
\renewcommand*{\tableautorefname}{Tabel}%
\renewcommand*{\partautorefname}{Bagian}%
\renewcommand*{\appendixautorefname}{Lampiran}%
\renewcommand*{\chapterautorefname}{Bab}%
\renewcommand*{\sectionautorefname}{Subbab}%
\renewcommand*{\subsectionautorefname}{Subsubbab}%
\renewcommand*{\subsubsectionautorefname}{Subsubsubbab}%
\renewcommand*{\paragraphautorefname}{paragraf}%
\renewcommand*{\subparagraphautorefname}{subparagraf}%
\renewcommand*{\FancyVerbLineautorefname}{garis}%
\renewcommand*{\theoremautorefname}{Teorema}%
\renewcommand*{\pageautorefname}{halaman}%

% Format to ignore underflow hbadness
\hbadness=99999
%--------------------------------------------------------------------%
%
% Hypenation untuk Bahasa Indonesia
%
% @author Petra Barus
%
%--------------------------------------------------------------------%
%
% Secara otomatis LaTeX dapat langsung memenggal kata dalam dokumen,
% tapi sering kali terdapat kesalahan dalam pemenggalan kata. Untuk
% memperbaiki kesalahan pemenggalan kata tertentu, cara pemenggalan
% kata tersebut dapat ditambahkan pada dokumen ini. Pemenggalan
% dilakukan dengan menambahkan karakter '-' pada suku kata yang
% perlu dipisahkan.
%
% Contoh pemenggalan kata 'analisa' dilakukan dengan 'a-na-li-sa'
%
%--------------------------------------------------------------------%

\hyphenation {
  % A
  %
  a-na-li-sa
  a-pli-ka-si
  a-lo-ka-si
  an-ta-ra
  ada-nya
  akan
  % B
  %
  be-be-ra-pa
  ber-ge-rak
  be-ri-kut
  ber-ko-mu-ni-ka-si
  buah
  Bouvet
  ber-fo-kus
  ber-fung-si
  ber-ja-lan
  % C
  %
  ca-ri
  con-strained
  Carzaniga
  cloud
  CloudFormation
  con-tai-ner
  ClusterIP
  % D
  %
  da-e-rah
  di-nya-ta-kan
  de-fi-ni-si
  di-bu-tuh-kan
  di-gu-na-kan
  di-tam-bah-kan-nya
  di-tem-pat-kan
  di-la-ku-kan
  de-ploy-ment
  di-kem-bang-kan
  di-im-ple-men-ta-si-kan
  de-activation
  da-pat
  di-ka-te-go-ri-kan
  de-ngan
  di-kem-bang-kan
  da-ta
  % E
  %
  e-ner-gi
  eks-klu-sif
  eks-ter-nal
  % F
  %
  fa-si-li-tas
  % G
  %
  ga-bung-an
  % H
  %
  ha-lang-an
  ha-sil
  hell
  % I
  % 
  i-nduk
  in-for-ma-si
  im-ple-men-tasi
  % J
  %
  % K
  %
  kom-po-si-si
  ka-re-na
  ke-sa-ba-ran-nya
  ka-me-ra
  kua-li-tas
  ke-na-ngan
  kom-plek-si-tas
  ke-ti-ka
  ke-le-bi-han-nya
  ke-gi-a-tan
  ko-mu-ni-ka-si
  ke-cil
  % L
  %
  la-ya-nan
  % M
  %
  me-ngu-ra-ngi
  meng-eva-lu-a-si
  me-nge-lo-la
  men-da-lam
  men-ja-lan-kan
  mak-si-mal
  me-nye-le-sai-kan
  me-ngunjungi
  men-du-kung
  me-nu-rut
  me-la-ku-kan
  mem-buat
  men-daftar-kan
  me-ngu-sul-kan
  me-mi-li-ki
  meng-gu-na-kan
  men-ja-di
  me-ru-pa-kan
  men-ja-ga
  me-mu-dah-kan
  me-ne-rus-
  mem-pro-ses
  % N
  %
  Na-mun
  % O
  %
  ob-so-lete
  or-kes-tra-si
  oto-ma-ti-sa-si
  % P
  %
  pro-vi-der
  pe-ru-sa-ha-an
  pe-rang-kat
  pro-ses
  plat-form
  pro-duk-si
  pe-ne-li-tian
  pe-ru-ba-han
  pa-ra-dig-ma
  pe-man-tau-an
  pe-ngum-pu-lan
  pa-ckage
  % Q
  %
  quality
  % R
  %
  re-source
  re-mote
  % S
  se-la-in
  stan-dar-di-sasi
  se-cu-ri-ty
  so-lu-si
  se-lu-ruh
  Soft-ware
  soft-ware
  se-buah
  se-ca-ra
  %
  % T
  % 
  ter-li-bat
  ter-pi-sah-kan
  ter-mi-nal
  ter-ba-tas
  % U
  %
  un-tuk
  % V
  %
  % W
  %
  % X
  %
  % Y
  % 
  % Z
  %
  zoo-keeper
}

%--------------------------------------------------------------------%
%
% Hypenation for English
%
% @author Muhammad Garebaldhie
%
%--------------------------------------------------------------------%
%
% Secara otomatis LaTeX dapat langsung memenggal kata dalam dokumen,
% tapi sering kali terdapat kesalahan dalam pemenggalan kata. Untuk
% memperbaiki kesalahan pemenggalan kata tertentu, cara pemenggalan
% kata tersebut dapat ditambahkan pada dokumen ini. Pemenggalan
% dilakukan dengan menambahkan karakter '-' pada suku kata yang
% perlu dipisahkan.
%
% Contoh pemenggalan kata 'analisa' dilakukan dengan 'a-na-li-sa'
%
%--------------------------------------------------------------------%


\makeatletter

\makeatother

\addbibresource{references.example.bib}

\begin{document}

\title{Judul TA}
\date{}
\author{
  Muhammad Garebaldhie ER Rahman \\
  NIM: 13520029
}
\newcommand\tanggalpengesahan{7 Juli 2024}

\pagenumbering{roman}
\setcounter{page}{1}

\clearpage
\pagestyle{empty}

\begin{center}
  \smallskip
  
  \Large \bfseries \MakeUppercase{\thetitle}
  \vfill
  
  \Large Laporan Tugas Akhir
  \vfill
  
  \large Disusun sebagai syarat kelulusan tingkat sarjana
  \vfill
  
  \large Oleh
  
  \Large \theauthor
  
  \vfill
  \begin{figure}[ht]
    \centering
    \includegraphics[width=0.15\textwidth]{cover-ganesha.jpg}
  \end{figure}
  \vfill
  
  \large
  \uppercase{
    Program Studi Teknik Informatika \\
    Sekolah Teknik Elektro \& Informatika \\
    Institut Teknologi Bandung
  }
  
  Juli 2024
  
\end{center}

\clearpage

\clearpage
\pagestyle{empty}

\begin{center}
  \smallskip
  
  \Large \bfseries \MakeUppercase{\thetitle}
  \vfill
  
  \Large Laporan Tugas Akhir
  \vfill
  
  \large Oleh
  
  \Large \theauthor
  
  \large Program Studi Teknik Informatika \\
  
  \normalsize \normalfont
  Sekolah Teknik Elektro dan Informatika \\
  Institut Teknologi Bandung \\
  
  \vfill
  \normalsize \normalfont
  Telah disetujui dan disahkan sebagai Laporan Tugas Akhir \\
  % Telah disetujui dan disahkan sebagai Draft Laporan Tugas Akhir \\
  di Bandung, pada tanggal \tanggalpengesahan
  
  \vspace{0.5cm}
  Pembimbing,
  
  \vfill
  \underline{Dr. techn. Muhammad Zuhri Catur Candra, S.T, M.T.
  } \\
  NIP. 19770921 201012 1 002
  
\end{center}
\clearpage

\chapter*{Lembar Pernyataan}

Dengan ini saya menyatakan bahwa:

\begin{enumerate}

	\item Pengerjaan dan penulisan Laporan Tugas Akhir ini dilakukan tanpa menggunakan bantuan yang tidak dibenarkan.
	\item Segala bentuk kutipan dan acuan terhadap tulisan orang lain yang digunakan di dalam penyusunan laporan tugas akhir ini telah dituliskan dengan baik dan benar.
	\item Laporan Tugas Akhir ini belum pernah diajukan pada program pendidikan di perguruan tinggi mana pun.
	      
\end{enumerate}

Jika terbukti melanggar hal-hal di atas, saya bersedia dikenakan sanksi sesuai dengan Peraturan Akademik dan Kemahasiswaan Institut Teknologi Bandung bagian Penegakan Norma Akademik dan Kemahasiswaan khususnya Pasal 2.1 dan Pasal 2.2.
\vspace{15mm}

Bandung, \tanggalpengesahan

\vspace{1.5cm}
Muhammad Garebaldie ER Rahman \\
NIM 13520029


\pagestyle{plain}

\clearpage
\chapter*{ABSTRAK}
\addcontentsline{toc}{chapter}{ABSTRAK}
\begin{center}
	\center
	\begin{singlespace}
		\large\bfseries\MakeUppercase{\thetitle}
		
		\normalfont\normalsize
		Oleh:
		
		\bfseries \theauthor
	\end{singlespace}
\end{center}

\begin{singlespace}
	\small
	Lorem ipsum dolor sit amet . Operator grafis dan tipografi mengetahui hal ini dengan baik, pada kenyataannya semua profesi yang berhubungan dengan alam semesta komunikasi memiliki hubungan yang stabil dengan kata-kata ini, tetapi apa itu? Lorem ipsum adalah teks dummy tanpa arti. Ini adalah urutan kata Latin yang, sebagaimana posisinya, tidak membentuk kalimat dengan pengertian yang utuh, tetapi memberikan kehidupan pada teks uji yang berguna untuk mengisi ruang-ruang yang selanjutnya akan ditempati dari teks ad hoc yang disusun oleh para profesional komunikasi.
	
	Ini tentu saja yang paling terkenal teks placeholder bahkan jika ada versi berbeda yang dapat dibedakan dari urutan pengulangan kata-kata Latin. Lorem ipsum berisi tipografi yang lebih banyak digunakan, sebuah aspek yang memungkinkan Anda untuk memiliki gambaran umum tentang rendering teks dalam hal pilihan font dan d ukuran font. Jika mengacu pada Lorem ipsum, ekspresi yang digunakan berbeda, yaitu fill text , fiktif text , blind text atau placeholder text : singkatnya, artinya juga bisa nol, tetapi kegunaannya sangat jelas untuk bertahan selama berabad-abad dan menolak versi ironis dan modern yang datang dengan kedatangan web.
	
	\textbf{\textit{Kata kunci: Lorem, Ipsum, Lorem Ipsum }}
	
\end{singlespace}
\clearpage
% \clearpage
\chapter*{ABSTRACT}
\addcontentsline{toc}{chapter}{ABSTRACT}

\begin{center}
  \center
  \begin{singlespace}
    \large\bfseries\MakeUppercase{Judul TA Disini}
    
    \normalfont\normalsize
    By:
    
    \bfseries \theauthor
  \end{singlespace}
\end{center}


\begin{singlespace}
  \small
  Lorem Ipsum is simply dummy text of the printing and typesetting industry. Lorem Ipsum has been the industry's standard dummy text ever since the 1500s, when an unknown printer took a galley of type and scrambled it to make a type specimen book. It has survived not only five centuries, but also the leap into electronic typesetting, remaining essentially unchanged. It was popularised in the 1960s with the release of Letraset sheets containing Lorem Ipsum passages, and more recently with desktop publishing software like Aldus PageMaker including versions of Lorem Ipsum.
  
  It is a long established fact that a reader will be distracted by the readable content of a page when looking at its layout. The point of using Lorem Ipsum is that it has a more-or-less normal distribution of letters, as opposed to using 'Content here, content here', making it look like readable English. Many desktop publishing packages and web page editors now use Lorem Ipsum as their default model text, and a search for 'lorem ipsum' will uncover many web sites still in their infancy. Various versions have evolved over the years, sometimes by accident, sometimes on purpose (injected humour and the like).
  
  Contrary to popular belief, Lorem Ipsum is not simply random text. It has roots in a piece of classical Latin literature from 45 BC, making it over 2000 years old. Richard McClintock, a Latin professor at Hampden-Sydney College in Virginia, looked up one of the more obscure Latin words, consectetur, from a Lorem Ipsum passage, and going through the cites of the word in classical literature, discovered the undoubtable source. Lorem Ipsum comes from sections 1.10.32 and 1.10.33 of "de Finibus Bonorum et Malorum" (The Extremes of Good and Evil) by Cicero, written in 45 BC. This book is a treatise on the theory of ethics, very popular during the Renaissance. The first line of Lorem Ipsum, "Lorem ipsum dolor sit amet..", comes from a line in section 1.10.32.
  
  \textbf{\textit{Keywords: Lorem, Ipsum, Lorem Ipsum }}
\end{singlespace}
\clearpage

\clearpage

\chapter*{Kata Pengantar}
\addcontentsline{toc}{chapter}{KATA PENGANTAR}

Puji dan syukur penulis panjatkan kepada Tuhan Yang Maha Esa atas berkat dan rahmatnya, laporan tugas akhir yang berjudul "\thetitle" dapat diselesaikan dalam rangka memenuhi syarat kelulusan tingkat sarjana. Perlu diakui pengerjaan tugas akhir ini didukung oleh banyak pihak. Khususnya, penulis ingin mengucapkan terima kasih kepada:

\begin{enumerate}
	\item Bapak Dr.techn. Muhammad Zuhri Catur Candra, S.T., M.T., selaku dosen pembimbing atas segala bentuk dukungan yang telah diberikan dan kesabarannya dalam membimbing penulis serta memberikan saran dalam pengerjaan tugas akhir.
	\item Bapak Yudistira Dwi Wardhana Asnar, S.T, Ph.D dan Dr. Agung Dewandaru, S.T., M.Sc., selaku dosen penguji atas segala masukan dan kritik yang telah diberikan terhadap tugas akhir penulis.
	\item Dicky Prima Satya, S.T, M.T., Bapak Adi Mulyanto, S.T, M.T., Robithoh Annur, S.T., M.Eng., Ph.D., dan Tricya Esterina Widagdo, ST., M.Sc. selaku dosen koordinator tim tugas akhir atas usahanya mengingatkan mahasiswa program studi Teknik Informatika untuk mengerjakan tugas akhirnya.
	\item Seluruh dosen program studi Teknik Informatika ITB yang telah memberikan ilmu pengetahuan yang sangat berharga bagi penulis.
	\item Ibu Rini Liani dan Bapak Edhie Hikmat selaku kedua orangtua penulis atas dukungan yang diberikan
	\item Rumah Amal Salman ITB serta tim Beasiswa Perintis yang telah membantu penulis sehingga penulis dapat menempuh pendidikan di Institut Teknologi Bandung dengan mudah.
	\item Teman-teman INIT 2020 yang telah menemani, memberikan inspirasi, serta dukungan moral kepada penulis dalam menempuh kuliah pada program studi Teknik Informatika.
	\item Teman-teman penulis khususnya anggota dari grup "temenin ngerjain TA", "Koordinasi penonton sempro", "Para Ajudan Pecinta Sedekah", "Kos Aufa Enjoyer", "kaliMANTAN", serta "karimun" yang telah memberikan kenangan berharga, motivasi, hiburan, serta bantuan untuk segala situasi.
	\item Sahabat Penulis khususnya Erik dan Rachel yang telah pantang menyerah berjuang bersama dalam berkompetisi CTF.
	\item Sahabat terdekat penulis khususnya Marcho, Aira, Gagas, Rio, Dhika, Kinan, Anca, Dipa, Ubai, Azka, Sarah, Dea, Fay, Epi, Aufa, Syahrul, Rifqi dan Malik yang telah menemani perjuangan dari TPB hingga saat ini, menjadi \textit{emotional support} di segala situasi, membantu penulis dalam proses pengejaan tugas akhir, serta membuat hari - hari menjadi lebih berwarna.
	\item Seluruh pihak lain yang tidak bisa disebutkan disini yang telah membantu dalam proses pengerjaan tugas akhir.
\end{enumerate}

Akhir kata, penulis mengucapkan terima kasih kepada semua pihak yang telah terlibat dalam pengerjaan tugas akhir ini. Penulis juga ingin menyampaikan mohon maaf apabila terdapat kesalahan maupun kekurangan dalam laporan tugas akhir ini. Penulis berharap semoga tugas akhir ini dapat bermanfaat bagi pembaca dan riset-riset kedepannya.

\begin{flushright}
	\vspace{0.5cm}
	Bandung, \tanggalpengesahan
	
	
	\vspace{1.5cm}
	
	Muhammad Garebaldhie ER Rahman
\end{flushright}

\titleformat*{\section}{\centering\bfseries\Large\MakeUpperCase}
\titlespacing*{\chapter}{0pt}{0pt}{4pt}

% Setting judul toc, lot, lof, bib
\renewcommand{\contentsname}{DAFTAR ISI}
\renewcommand{\listfigurename}{DAFTAR GAMBAR}
\renewcommand{\listtablename}{DAFTAR TABEL}
\renewcommand{\bibname}{DAFTAR PUSTAKA}

% daftar isi, lampiran, gambar, table
\tableofcontents
\listofappendices
\listoffigures
\listoftables

\newpage

\titleformat*{\section}{\bfseries\large}
\pagenumbering{arabic}

%----------------------------------------------------------------%
% Konfigurasi Bab
%----------------------------------------------------------------%
\setcounter{page}{1}
\renewcommand{\chaptername}{BAB}
\renewcommand{\thechapter}{\Roman{chapter}}
%----------------------------------------------------------------%

%----------------------------------------------------------------%
% Dafter Bab
% Untuk menambahkan daftar bab, buat berkas bab misalnya `chapter-6` di direktori `chapters`, dan masukkan ke sini.
%----------------------------------------------------------------%
\chapter{Pendahuluan}

Bab Pendahuluan secara umum yang dijadikan landasan kerja dan arah kerja penulis tugas akhir, berfungsi mengantar pembaca untuk membaca laporan tugas akhir secara keseluruhan.

\section{Latar Belakang}
\label{sec:latarbelakang}

Latar Belakang berisi dasar pemikiran, kebutuhan atau alasan yang menjadi ide dari topik tugas akhir. Tujuan utamanya adalah untuk memberikan informasi secukupnya kepada pembaca agar memahami topik yang akan dibahas.  Saat menuliskan bagian ini, posisikan anda sebagai pembaca – apakah anda tertarik untuk terus membaca?

\section{Rumusan Masalah}

Rumusan Masalah berisi masalah utama yang dibahas dalam tugas akhir. Rumusan masalah yang baik memiliki struktur sebagai berikut:

\begin{enumerate}
	\item Penjelasan ringkas tentang kondisi/situasi yang ada sekarang terkait dengan topik utama yang dibahas Tugas Akhir.
	\item Pokok persoalan dari kondisi/situasi yang ada, dapat dilihat dari kelemahan atau kekurangannya. \textbf{Bagian ini merupakan inti dari rumusan masalah}.
	\item Elaborasi lebih lanjut yang menekankan pentingnya untuk menyelesaikan pokok persoalan tersebut.
	\item Usulan singkat terkait dengan solusi yang ditawarkan untuk menyelesaikan persoalan.
\end{enumerate}

Penting untuk diperhatikan bahwa persoalan yang dideskripsikan pada subbab ini akan dipertanggungjawabkan di bab Evaluasi apakah terselesaikan atau tidak.

\section{Tujuan}

Tuliskan tujuan utama dan/atau tujuan detil yang akan dicapai dalam pelaksanaan tugas akhir. Fokuskan pada hasil akhir yang ingin diperoleh setelah tugas akhir diselesaikan, terkait dengan penyelesaian persoalan pada rumusan masalah. Penting untuk diperhatikan bahwa tujuan yang dideskripsikan pada subbab ini akan dipertanggungjawabkan di akhir pelaksanaan tugas akhir apakah tercapai atau tidak.

\section{Batasan Masalah}

Tuliskan batasan-batasan yang diambil dalam pelaksanaan tugas akhir. Batasan ini dapat dihindari (tidak perlu ada) jika topik/judul tugas akhir dibuat cukup spesifik.

\section{Metodologi}

Tuliskan semua tahapan yang akan dilalui selama pelaksanaan tugas akhir. Tahapan ini spesifik untuk menyelesaikan persoalan tugas akhir. Tahapan studi literatur tidak perlu dituliskan karena ini adalah pekerjaan yang harus Anda lakukan selama proses pelaksanaan tugas akhir.

\section{Jadwal Pelaksanaan Tugas Akhir}

Tuliskan rencana kegiatan dan jadwal (dirinci sampai per minggu) mulai dari awal pelaksanaan Tugas Akhir I s.d. sidang tugas akhir berikut milestones dan deliverables yang harus diberikan. Jadwal ini dapat dibantu dengan membuat sebuah tabel timeline.
\chapter{Studi Literatur}
\label{chapter:studi-literatur}

Bab Studi Literatur digunakan untuk mendeskripsikan kajian literatur yang terkait dengan persoalan tugas akhir. Tujuan studi literatur adalah:

\begin{enumerate}
  \item menunjukkan kepada pembaca adanya gap seperti pada rumusan masalah yang memang belum terselesaikan,
  \item memberikan pemahaman secukupnya kepada pembaca tentang teori atau pekerjaan terkait yang terkait langsung dengan penyelesaian persoalan, serta
  \item menyampaikan informasi apa saja yang sudah ditulis/dilaporkan oleh pihak lain (peneliti/Tugas Akhir/Tesis) tentang hasil penelitian/pekerjaan mereka yang sama atau mirip kaitannya dengan persoalan tugas akhir.
\end{enumerate}

Kita juga dapat memisah beberapa bagian latex untuk \textit{readability}. Kita dapat memasukan file latex lainnya dengan menggunakan fitur input. Berikut merupakan contoh ketika memasukan bagian contoh-subbab ke chapter-2

\section{Contoh Subbab}
Perujukan literatur dapat dilakukan dengan menambahkan entri baru di berkas. Tulisan ini merujuk pada \parencite{abugessaisa2023} atau \parencite{appliedpredictivemodel} dan \parencite{vogler2015diane}

Sekarang mau ke bab berapa yaaaa.... hmm... ke bab \ref{chapter:studi-literatur} ahhhhh.

\blindtext

\subsection{Contoh Subsubbab}
\label{subsec:contoh-subsec}

\blindtext

\begin{figure}[ht]
	\centering
	\includegraphics[width=0.5\textwidth]{resources/cover-ganesha.jpg}
	\caption{Contoh gambar}
\end{figure}

\subsubsection{Subsubsubbab}
\blindtext

\subsubsubsection{sub sub sub sub bab}
\blindtext

\begin{table}[h]
	\caption{Tabel random}
	\vspace{0.25cm}
	\begin{center}
		\begin{tabular}{|c|c|c|c|}
			\hline
			Title1 & Title2 & Title3 & Title4  \tabularnewline
			\hline
			1647   & 1.97   & 0.68   & 1.90 \tabularnewline
			2301   & 2.92   & 1.06   & 2.75 \tabularnewline
			2969   & 3.23   & 1.16   & 3.78 \tabularnewline
			3791   & 4.39   & 1.40   & 4.14 \tabularnewline
			4625   & 6.72   & 1.87   & 5.59 \tabularnewline
			\hline
		\end{tabular}
	\end{center}
\end{table}


\section{Menyisipkan Persamaan}

Beberapa contoh menyisipkan persamaan.


\subsection{Contoh Bikin Equation}
\textbf{text tebal} dan ini \emph{miring}, bikin persamaan di baris yang sama, tinggal pake dolar2 $\Psi(\vec{r}_1,...,\vec{r}_N)$, sehingga persamaan Schr\"{o}dinger, terus, persamaan yang dinomeri kayak gini 
%ini contoh bikin persamaan, ..... :D
\begin{equation}
  \left[ \sum_{i}^{N}-\frac{\hbar^2}{2m}\nabla_i^2 + \sum_{i}^{N}V(\vec{r}_i)+ \sum_{i<j}^{N}(\vec{r}_i,\vec{r}_j)\right]\Psi = E\Psi 
\end{equation}

untuk $N$-elektron, dengan $\hat{H}$=Hamiltonian, $E$=Energi total, $\hat{T}$=Energi kinetik, $\hat{V}$=Energi potensial, dan $\hat{U}$=Interaksi ektron-elektron.

\subsection{Bikin Matrix}
Lalalallala.... bikin matrix sekarang, yang ini dikecilin, pake smaller
  {\smaller
    \begin{equation}
      \Psi({\bf r}_1, {\bf r}_2, \cdots {\bf r}_N) = \frac{1}{\sqrt{N!}}\left| \begin{array}{llcl}
        \phi_1({\bf r}_1)     & \phi_2({\bf r}_1)     & \cdots                & \phi_N({\bf r}_1)     \\
        \phi_1({\bf r}_2)     & \phi_2({\bf r}_2)     & \cdots                & \phi_N({\bf r}_2)     \\
        \phi_1({\bf r}_3)     & \phi_2({\bf r}_3)     & \cdots                & \phi_N({\bf r}_3)     \\
        \multicolumn{1}{c}{.} & \multicolumn{1}{c}{.} & \multicolumn{1}{c}{.} & \multicolumn{1}{c}{.} \\
        \multicolumn{1}{c}{.} & \multicolumn{1}{c}{.} & \multicolumn{1}{c}{.} & \multicolumn{1}{c}{.} \\
        \multicolumn{1}{c}{.} & \multicolumn{1}{c}{.} & \multicolumn{1}{c}{.} & \multicolumn{1}{c}{.} \\
        \phi_1({\bf r}_N)     & \phi_2({\bf r}_N)     & \cdots                & \phi_N({\bf r}_N)     \\
      \end{array} \right|
    \end{equation}
  }
% \blankpage
\chapter{Analisis Persoalan dan Rancangan Solusi}

Tujuan utama penulisan bab ini adalah untuk menguraikan rencana penyelesaian masalah implementasi dari judul TA


\section{Analisis}
\blindtext

\section{Rancangan}
\blindtext

\chapter{Implementasi dan Pengujian}
Bab ini akan menjelaskan proses implementasi dari rancangan solusi yang telah dikaji pada Bab III. Setelah pembahasan terkait implementasi, akan dilanjutkan dengan pemaparan hasil uji terkait implementasi yang telah dibuat.

\section{Lingkungan}
\blindtext

\section{Implementasi}
\blindtext

\section{Pengujian}
\blindtext



\chapter{Penutup}

\section{Kesimpulan}
\blindtext

\section{Saran}
\blindtext
%---------------------------------------------------------------%

% Daftar pustaka
\printbibliography

% Setting judul lampiran
\titlespacing*{\chapter}{0pt}{0pt}{0pt}
\titlespacing*{\section}{0pt}{0pt}{*1}

% Setting judul anak lampiran
\titleformat*{\section}{\bfseries}

\appendix

\chapter{Contoh gambar pengujian}

\begin{figure}[ht]
  \centering
  \includegraphics[width=0.5\textwidth]{resources/cover-ganesha.jpg}
  \caption{Contoh gambar}
\end{figure}


\end{document}
